\documentclass[a4paper,12pt]{article}
\usepackage[utf8]{inputenc}
\usepackage[T1]{fontenc}
\usepackage[ngerman]{babel}
\usepackage{geometry}
\geometry{margin=2.5cm}
\usepackage{amsmath}
\usepackage{booktabs}
\usepackage{hyperref}
\hypersetup{colorlinks=true, linkcolor=blue, urlcolor=blue}

\begin{document}

% =========================================================
\section*{Batch Normalization: Accelerating Deep Network Training
          by Reducing Internal Covariate Shift}
% =========================================================

\noindent\textbf{Autoren:} Sergey Ioffe, Christian Szegedy\\[2pt]
\textbf{Institution:} Google Inc.\\[2pt]
\textbf{Ver\"offentlicht:} arXiv:1502.03167v3 [cs.LG], 2015 (ICML 2015)\\[2pt]

% =========================================================
\subsection*{Motivation und Problemstellung}
% =========================================================

Das Training tiefer neuronaler Netze wird dadurch erschwert, dass sich die
Eingangsverteilung jeder Schicht w\"ahrend des Trainings kontinuierlich ver\"andert,
da sich die Parameter der vorgelagerten Schichten laufend anpassen. Dieses Ph\"anomen
bezeichnen die Autoren als \emph{Internal Covariate Shift}. Es erzwingt niedrige
Lernraten und eine sorgf\"altige Initialisierung, verlangsamt die Konvergenz
erheblich und macht das Training von Netzen mit s\"attigenden Nichtlinearit\"aten
(z.\,B.\ Sigmoid) besonders schwierig.

Als L\"osung schlagen die Autoren \textbf{Batch Normalization (BN)} vor: eine
Normierungsschicht, die die Eingangsverteilung jeder Schicht pro Mini-Batch
stabilisiert und direkt in die Netzwerkarchitektur eingebettet wird.

% =========================================================
\subsection*{Kernbeitrag: Batch Normalizing Transform}
% =========================================================

F\"ur eine Aktivierung $x$ \"uber einen Mini-Batch $\mathcal{B} = \{x_1,\ldots,x_m\}$
wird die BN-Transformation definiert als:
\[
  \hat{x}_i = \frac{x_i - \mu_\mathcal{B}}{\sqrt{\sigma^2_\mathcal{B} + \epsilon}},
  \qquad
  y_i = \gamma\,\hat{x}_i + \beta,
\]
wobei $\mu_\mathcal{B}$ und $\sigma^2_\mathcal{B}$ Mini-Batch-Mittelwert und
-Varianz sind, $\epsilon$ eine numerische Konstante, und $\gamma, \beta$ gelernte
Skalierungs- und Verschiebungsparameter. Diese Parameter stellen sicher, dass die
BN-Schicht die Identit\"atstransformation repr\"asentieren kann und die
Ausdrucksst\"arke des Netzes erhalten bleibt.

F\"ur faltende Schichten wird BN auf alle Aktivierungen einer Feature-Map
\"uber Mini-Batch und r\"aumliche Positionen gemeinsam angewendet;
dabei werden $\gamma$ und $\beta$ pro Feature-Map (nicht pro Aktivierung) gelernt.

W\"ahrend der Inferenz werden Populations- statt Mini-Batch-Statistiken verwendet,
so dass die Normierung eine einfache lineare Transformation darstellt, die
mit den gelernten $\gamma, \beta$ zusammengefasst werden kann.

% =========================================================
\subsection*{Vorteile von Batch Normalization}
% =========================================================

\begin{itemize}
  \item \textbf{H\"ohere Lernraten:} BN begrenzt die Amplifikation von
    Parameteraktualisierungen durch tiefe Schichten und reduziert das Risiko
    des Gradientenexplodierens oder -verschwindens.

  \item \textbf{Reduzierter Regularisierungsbedarf:} Durch die Abh\"angigkeit
    jedes Beispiels von den \"ubrigen Beispielen im Mini-Batch wirkt BN
    implizit regularisierend; Dropout kann in batch-normalisierten Netzen
    oft reduziert oder weggelassen werden.

  \item \textbf{Robustheit gegen\"uber Initialisierung:} Da BN die
    Eingangsverteilung stabilisiert, sind Netze weniger empfindlich gegen\"uber
    der genauen Wahl der Startgewichte.
\end{itemize}

% =========================================================
\subsection*{Experimentelle Ergebnisse}
% =========================================================

Die Autoren wenden BN auf eine Variante des Inception-Netzes (ImageNet-Klassifikation)
an und zeigen:

\begin{itemize}
  \item \textbf{Beschleunigung:} BN-Inception (BN-x5) erreicht die Genauigkeit
    des Basis-Inception-Netzes in 14-mal weniger Trainingsschritten.

  \item \textbf{H\"ohere Endgenauigkeit:} BN-x30 erzielt nach $6\cdot10^6$
    Schritten 74{,}8\,\% Top-1-Genauigkeit, gegen\"uber 72{,}2\,\% des
    Original-Inception nach $31\cdot10^6$ Schritten.

  \item \textbf{State of the Art (ImageNet 2015):} Ein Ensemble batch-normalisierter
    Netze erreicht 4{,}9\,\% Top-5-Fehler auf dem Validierungsdatensatz und
    4{,}82\,\% auf dem Testdatensatz -- besser als alle damals bekannten
    Systeme und erstmals vergleichbar mit menschlicher Leistung.
\end{itemize}

% =========================================================
\subsection*{Bezug zur Masterarbeit}
% =========================================================

\subsubsection*{1.\ Einsatz im U-Net-Surrogatmodell}

Batch Normalization wird im U-Net der Masterarbeit nach jeder $3\times3$-Faltung
und vor der ReLU-Aktivierung eingesetzt (sowohl im Encoder als auch im Decoder).
Dies entspricht der von Ioffe \& Szegedy empfohlenen Platzierung \emph{vor} der
Nichtlinearit\"at. BN stabilisiert das Training bei variablen Kristallkonfigurationen,
bei denen Eingangsskalen und Feature-Verteilungen stark schwanken k\"onnen.

\subsubsection*{2.\ Kein Einsatz im FNO}

Das FNO-Modell dieser Arbeit verwendet kein Batch Normalization. Stattdessen
\"ubernimmt der spektrale Gl\"attungsbias durch Mode-Abschneidung eine implizit
regularisierende Funktion, und GELU-Aktivierungen ohne explizite Normierung
haben sich in der Operator-Learning-Literatur als stabil erwiesen. Der Verzicht
auf BN im FNO ist architekturmotiviert und im Einklang mit der Originalimplementierung
von \textcite*{li_fourier_2021}.

\subsubsection*{3.\ Dropout-Verzicht}

Konsistent mit der Empfehlung von Ioffe \& Szegedy, dass BN Dropout weitgehend
ersetzen kann, wird im U-Net kein Dropout verwendet. Die Regularisierung erfolgt
stattdessen durch geometrische Vielfalt in den Trainingsdaten.

% =========================================================
\subsection*{Kritische Einsch\"atzung}
% =========================================================

\begin{itemize}
  \item Sp\"atere Arbeiten (z.\,B.\ Ba et al., 2016: Layer Normalization;
    He et al., 2016: Group Normalization) haben gezeigt, dass BN f\"ur kleine
    Mini-Batches und wiederkehrende Netze problematisch sein kann, da die
    Mini-Batch-Statistiken dann instabil werden. Im vorliegenden Kontext
    (Batch-Gr\"o\ss en 8 und 16, konvolutionelles U-Net) ist dies kein relevantes
    Problem.

  \item Die theoretische Erkl\"arung \"uber Internal Covariate Shift wurde sp\"ater
    in Frage gestellt (Santurkar et al., 2018: \emph{How does BN help optimization?});
    empirisch ist der Nutzen von BN jedoch unbestritten.
\end{itemize}

\end{document}
